\subsection{Worked example: translating a bounded response}
\label{subsec:worked-translation}
Consider
\[
  f=(\Eventually_{\le5}p)\land(\Always_{<2}\neg p).
\]
The property requires the first occurrence of $p$ to be no earlier than two
and no later than five time units after the temporal origin. Figure~\ref{fig:translation-reference-tba}
shows a small TBA for this language. The construction below follows the
formula through the clocked-LTL intermediate language and the timed
reinterpretation.

\begin{figure}[t]
\centering
\resizebox{0.88\columnwidth}{!}{\begin{tikzpicture}[node distance=22mm]
  \node[fmstate] (q0) {$q_{\mathit{init}}$};
  \node[fmstate,right=of q0] (q1) {$q_{\mathit{wait}}$};
  \node[fmaccept,right=of q1] (q2) {$q_{\mathit{acc}}$};
  \draw[fmedge] ([xshift=-8mm]q0.west) -- (q0.west);
  \draw[fmedge] (q0) -- node[fmlabel,above] {$\neg p\;/\;\mathsf{rst}(x)$} (q1);
  \draw[fmedge] (q1) edge[loop above] node[fmlabel] {$x\le5,\;\neg p$} (q1);
  \draw[fmedge] (q1) -- node[fmlabel,above] {$2\le x\le5,\;p$} (q2);
  \draw[fmedge] (q2) edge[loop above] node[fmlabel] {$\top$} (q2);
\end{tikzpicture}
}
\caption{A two-location TBA for $f$. The first location enforces the deadline
$x\le5$ and the guard $x\ge2$ requires $p$ only in the interval $[2,5]$.}
\label{fig:translation-reference-tba}
\end{figure}

\begin{enumerate}[label=Step~\arabic*:,leftmargin=*]
\item The formula is in negation normal form. Expanding the derived operators
and applying the timed-operator identities gives
\[
\begin{aligned}
f
 &= (\top\Until_{\le5}p)\land(\bot\Release_{<2}\neg p)\\
 &= \bigl(p\lor(\top\land\top\hUntil_{\le5}p)\bigr)\\
 &\qquad\land\bigl(\neg p\land(\bot\lor\bot\hRelease_{<2}\neg p)\bigr).
\end{aligned}
\]
The current-position alternatives in the last line are the derived-operator
cases of Eqs.~\eqref{eq:derive-ule} and~\eqref{eq:derive-rle}.

\item The two distinct primitive timed subformulas receive clocks $x$ and
$y$, respectively. Their clocked-LTL translation is
\begin{align*}
\Phi_U={}&p\lor\bigl(\top\land\Next\Bigl(
  ((x\le5)\land\Unch(x)\land\top)\Until\\[-1mm]
&\qquad((x\le5)\land p)\Bigr)\bigr),\\[-1mm]
\Phi_R={}&\neg p\land\bigl(\bot\lor(\Reset(y)\land\\[-1mm]
&\qquad\Next(\neg p\WeakUntil((y\ge2)\lor(\bot\land\neg p))))\bigr),\\[-1mm]
\T(f)={}&\Phi_U\land\Phi_R.
\end{align*}
The upper-bounded Until uses $x$ as a preservation clock. The lower test
$y\ge2$ is the complement of the strict source bound $<2$; the Release clause
therefore permits $p$ beginning exactly at age two.

\item Spot translates $\T(f)$ into the propositional B\"uchi automaton on the
left of Fig.~\ref{fig:translation-steps45}. It has five states and eleven
cube-labeled transitions. Clock guards, reset markers, preservation markers,
and action literals are independent propositions at this stage. In
particular, the cubes containing both $p$ and $\neg p$ are propositionally
satisfiable to Spot.

\item Relaxation removes negative literals over extended clock atoms, and
reset completion assigns resets according to the clock class: $x$ is reset
unless the transition carries $\Unch(x)$, while $y$ is reset exactly when it
carries $\Reset(y)$. The right side of Fig.~\ref{fig:translation-steps45}
shows the resulting TBA transitions after the markers are removed. This is
the boundary of the translation theorem; subsequent optimization removes
the unsatisfiable action cubes and simplifies the automaton.
\end{enumerate}

\begin{figure*}[!t]
\centering
\resizebox{\textwidth}{!}{\begin{tikzpicture}[
  x=1cm,y=1cm,
  lab/.style={font=\small,align=center,fill=white,inner sep=1.2pt}
]
% Left: Spot 2.16 output; state IDs match examples/spot/worked_translation.dot.
\begin{scope}[xshift=0cm]
  \node[fmnote,font=\small\bfseries] at (2.4,6.0) {Spot output};
  \node[fmstate] (l0) at (0.0,4.8) {$2$};
  \node[fmstate] (l1) at (0.0,2.4) {$3$};
  \node[fmaccept] (l2) at (4.8,4.8) {$1$};
  \node[fmstate] (l3) at (4.8,2.4) {$4$};
  \node[fmaccept] (l4) at (2.4,0.0) {$0$};
  \draw[fmedge] (-1.0,4.8) -- (l0);
  \draw[fmedge] (l0) -- node[lab,left] {$!p\land\neg p\land\mathsf{rst}(y)$} (l1);
  \draw[fmedge] (l0) -- node[lab,above] {$p\land\neg p\land\mathsf{rst}(y)$} (l2);
  \draw[fmedge] (l1) edge[loop left,looseness=5] node[lab,left] {$\begin{array}{r}\mathsf{unch}(x)\land x\le5\\{}\land\,!p\land\neg p\\{}\land\,!(y\ge2)\end{array}$} (l1);
  \draw[fmedge] (l1) -- node[lab,pos=.42,sloped,above] {$\begin{array}{c}x\le5\land p\land\neg p\\{}\land\,!(y\ge2)\end{array}$} (l2);
  \draw[fmedge] (l1) -- node[lab,below] {$\begin{array}{c}\mathsf{unch}(x)\land x\le5\\{}\land\,!p\land y\ge2\end{array}$} (l3);
  \draw[fmedge] (l1) -- node[lab,left,pos=.45] {$\begin{array}{r}x\le5\land p\\{}\land y\ge2\end{array}$} (l4);
  \draw[fmedge] (l2) edge[loop right,looseness=5] node[lab,right] {$\begin{array}{l}\neg p\\{}\land\,!(y\ge2)\end{array}$} (l2);
  \draw[fmedge] (l2) to[out=-15,in=10,looseness=1.3] node[lab,right,pos=.3] {$y\ge2$} (l4);
  \draw[fmedge] (l3) edge[loop right,looseness=5] node[lab,right] {$\begin{array}{l}\mathsf{unch}(x)\\{}\land x\le5\land\,!p\end{array}$} (l3);
  \draw[fmedge] (l3) -- node[lab,right,pos=.4] {$x\le5\land p$} (l4);
  \draw[fmedge] (l4) edge[loop below,looseness=5] node[lab,below] {$\top$} (l4);
\end{scope}

% Right: relaxation and reset completion, with the same Spot state IDs.
\begin{scope}[xshift=9.6cm]
  \node[fmnote,font=\small\bfseries] at (2.4,6.0) {After relaxation and reset completion};
  \node[fmstate] (r0) at (0.0,4.8) {$2$};
  \node[fmstate] (r1) at (0.0,2.4) {$3$};
  \node[fmaccept] (r2) at (4.8,4.8) {$1$};
  \node[fmstate] (r3) at (4.8,2.4) {$4$};
  \node[fmaccept] (r4) at (2.4,0.0) {$0$};
  \draw[fmedge] (-1.0,4.8) -- (r0);
  \draw[fmedge] (r0) -- node[lab,left] {$\begin{array}{r}\neg p\\/\;\mathsf{rst}(x),\mathsf{rst}(y)\end{array}$} (r1);
  \draw[fmedge] (r0) -- node[lab,above] {$p\land\neg p\;/\;\mathsf{rst}(x),\mathsf{rst}(y)$} (r2);
  \draw[fmedge] (r1) edge[loop left,looseness=5] node[lab,left] {$x\le5,\;\neg p$} (r1);
  \draw[fmedge] (r1) -- node[lab,pos=.42,sloped,above] {$x\le5,\;p\land\neg p\;/\;\mathsf{rst}(x)$} (r2);
  \draw[fmedge] (r1) -- node[lab,below] {$x\le5\land y\ge2,\;\neg p$} (r3);
  \draw[fmedge] (r1) -- node[lab,left,pos=.45] {$\begin{array}{r}x\le5\land y\ge2,\;p\\/\;\mathsf{rst}(x)\end{array}$} (r4);
  \draw[fmedge] (r2) edge[loop right,looseness=5] node[lab,right] {$\begin{array}{l}\neg p\\/\;\mathsf{rst}(x)\end{array}$} (r2);
  \draw[fmedge] (r2) to[out=-15,in=10,looseness=1.3] node[lab,right,pos=.3] {$\begin{array}{l}y\ge2,\;\top\\/\;\mathsf{rst}(x)\end{array}$} (r4);
  \draw[fmedge] (r3) edge[loop right,looseness=5] node[lab,right] {$x\le5,\;\neg p$} (r3);
  \draw[fmedge] (r3) -- node[lab,right,pos=.4] {$\begin{array}{l}x\le5,\;p\\/\;\mathsf{rst}(x)\end{array}$} (r4);
  \draw[fmedge] (r4) edge[loop below,looseness=5] node[lab,below] {$\top\;/\;\mathsf{rst}(x)$} (r4);
\end{scope}
\end{tikzpicture}
}
\caption{Generated automata for the worked formula. Left: Spot's B\"uchi
automaton, before timed reinterpretation. Right: after relaxation and reset
completion. Here $!p$ is the negative literal of the extended atom $p$,
whereas $\neg p$ is an action proposition; their conjunction is retained
until the separate TBA optimization stage. Edges on the right are labeled
``guard, action / resets.''}
\label{fig:translation-steps45}
\end{figure*}
